\documentclass[11pt,a4paper]{article}
% Preâmbulo comum. Não duplicar pacotes nos módulos.
\usepackage[utf8]{inputenc}
\usepackage[T1]{fontenc}
\usepackage[brazil]{babel}
\usepackage{amsmath,amssymb,amsthm,mathtools}
\usepackage{geometry}
\geometry{a4paper,margin=2.3cm}
\usepackage{microtype}
\usepackage{enumitem}
\usepackage{graphicx}
\usepackage{xcolor}
\usepackage{booktabs}
\usepackage{hyperref}
\hypersetup{colorlinks=true,linkcolor=blue!40!black,urlcolor=blue!50!black}

\setcounter{tocdepth}{2}
\numberwithin{equation}{section}

\theoremstyle{plain}
\newtheorem{teorema}{Teorema}[section]
\newtheorem{proposicao}[teorema]{Proposição}
\newtheorem{lema}[teorema]{Lema}
\newtheorem{corolario}[teorema]{Corolário}

\theoremstyle{definition}
\newtheorem{definicao}[teorema]{Definição}
\newtheorem{exemplo}[teorema]{Exemplo}
\newtheorem{exercicio}{Exercício}[section]

\theoremstyle{remark}
\newtheorem{observacao}[teorema]{Observação}
\newtheorem{contraexemplo}[teorema]{Contraexemplo}
\newtheorem{intuicao}[teorema]{Intuição}

\newcommand{\R}{\mathbb{R}}
\newcommand{\N}{\mathbb{N}}
\newcommand{\Z}{\mathbb{Z}}
\newcommand{\Q}{\mathbb{Q}}
\DeclareRobustCommand{\eps}{\varepsilon}
\DeclareRobustCommand{\vet}[1]{\vec{#1}}
\DeclareRobustCommand{\dif}{\mathrm{d}}
\DeclareRobustCommand{\abs}[1]{\left\lvert #1 \right\rvert}
\DeclareRobustCommand{\norm}[1]{\left\lVert #1 \right\rVert}
\DeclareMathOperator{\grad}{grad}
\DeclareMathOperator{\Div}{div}
\DeclareMathOperator{\rot}{rot}

\newcommand{\base}[1]{\par\medskip\noindent\textbf{Base para relembrar.} #1\par\medskip}
\newcommand{\cuidado}[1]{\par\medskip\noindent\textbf{Quando a hipótese falha.} #1\par\medskip}

\newcommand{\figura}[2]{%
  \begin{figure}[htbp]
  \centering
  \IfFileExists{#1}{%
    \includegraphics[width=0.84\textwidth]{#1}%
  }{%
    \fbox{\parbox{0.8\textwidth}{\centering Figura ainda não gerada:\\ \texttt{#1}}}%
  }
  \caption{#2}
  \end{figure}%
}

\newenvironment{esqueleto}{%
  \par\medskip\noindent\textbf{Esqueleto da demonstração.}\par\nopagebreak
  \begin{enumerate}[leftmargin=1.4em]
}{\end{enumerate}\medskip}

\title{Cálculo III: Green, Stokes e Gauss}
\author{}
\date{}
\begin{document}
\maketitle
\tableofcontents

\section{Três versões do mesmo fato}

O teorema fundamental do cálculo diz que a integral da derivada num intervalo é a diferença da função na fronteira, que nesse caso são dois pontos. Green, Stokes e Gauss repetem a frase em dimensões maiores. O que se integra no domínio é um invariante infinitesimal --- uma derivada, o rotacional ou a divergência. O que se lê na fronteira é circulação ou fluxo. A orientação da fronteira faz parte do enunciado: sem ela o sinal sai errado.

\begin{center}
\begin{tabular}{lll}
\toprule
região & invariante integrado & leitura na fronteira \\
\midrule
intervalo & derivada & diferença nos extremos \\
domínio plano & $\partial_x Q-\partial_y P$ & circulação anti-horária \\
superfície orientada & $\nabla\times\vet{F}$ & circulação do bordo \\
sólido & $\nabla\cdot\vet{F}$ & fluxo da normal para fora \\
\bottomrule
\end{tabular}
\end{center}

Green é Stokes no plano: a superfície é chata, a normal é $\vet{k}$ e a componente do rotacional que sobra é $\partial_x Q-\partial_y P$. Gauss é a versão de divergência: no lugar da circulação do bordo entra o fluxo da casca. Os três medem, por dentro, o que a fronteira já mostra por fora.

\base{Teorema fundamental do cálculo numa variável. Teorema de Clairaut: parciais mistas contínuas comutam. Num retângulo, uma função contínua tem integrais iteradas iguais (Fubini). No plano, $x=r\cos\theta$, $y=r\sin\theta$ e $\dif A=r\,\dif r\,\dif\theta$.}

\section{Green no retângulo}

\begin{teorema}
\label{thm:green-ret}
Seja $D=[a,b]\times[c,d]$, com $a<b$ e $c<d$. Sejam $P$ e $Q$ de classe $C^1$ num aberto que contém $D$. Oriente $\partial D$ no sentido anti-horário: a base da esquerda para a direita, o lado direito de baixo para cima, o topo da direita para a esquerda, o lado esquerdo de cima para baixo. Então
\[
\oint_{\partial D} P\,\dif x+Q\,\dif y
=\iint_D\bigl(\partial_x Q-\partial_y P\bigr)\,\dif A.
\]
\end{teorema}

\begingroup\catcode`\_=12
\figura{../figuras/green.png}{Retângulo com a fronteira anti-horária.}
\endgroup

\begin{proof}
Nomeie os lados: $C_1$ a base, $C_2$ a direita, $C_3$ o topo, $C_4$ a esquerda.

No cálculo de $\oint P\,\dif x$ os lados verticais não contribuem, porque neles $x$ é constante e $\dif x=0$. Na base, $y=c$ e $x$ vai de $a$ até $b$. No topo, $y=d$ e $x$ vai de $b$ até $a$, o que troca o sinal:
\begin{align*}
\oint_{\partial D} P\,\dif x
&=\int_a^b P(x,c)\,\dif x+\int_b^a P(x,d)\,\dif x
=\int_a^b P(x,c)\,\dif x-\int_a^b P(x,d)\,\dif x.
\end{align*}
Para cada $x$ fixo, o teorema fundamental do cálculo na variável $y$ dá
\[
P(x,d)-P(x,c)=\int_c^d \partial_y P(x,y)\,\dif y,
\]
porque $\partial_y P$ existe e é contínua no retângulo. Portanto
\[
\oint_{\partial D} P\,\dif x
=\int_a^b\Biggl(-\int_c^d \partial_y P(x,y)\,\dif y\Biggr)\dif x
=-\int_a^b\int_c^d \partial_y P\,\dif y\,\dif x
=-\iint_D \partial_y P\,\dif A.
\]

No cálculo de $\oint Q\,\dif y$ os lados horizontais não contribuem, porque neles $\dif y=0$. No lado direito, $x=b$ e $y$ vai de $c$ até $d$. No lado esquerdo, $x=a$ e $y$ vai de $d$ até $c$:
\begin{align*}
\oint_{\partial D} Q\,\dif y
&=\int_c^d Q(b,y)\,\dif y+\int_d^c Q(a,y)\,\dif y
=\int_c^d\bigl(Q(b,y)-Q(a,y)\bigr)\,\dif y.
\end{align*}
Para cada $y$ fixo,
\[
Q(b,y)-Q(a,y)=\int_a^b \partial_x Q(x,y)\,\dif x,
\]
logo
\[
\oint_{\partial D} Q\,\dif y
=\int_c^d\int_a^b \partial_x Q\,\dif x\,\dif y
=\iint_D \partial_x Q\,\dif A.
\]
A troca da ordem é o teorema de Fubini: $\partial_x Q$ é contínua no compacto $D$. Somando as duas identidades,
\[
\oint_{\partial D} P\,\dif x+Q\,\dif y
=\iint_D\bigl(\partial_x Q-\partial_y P\bigr)\,\dif A.
\]
\end{proof}

\begin{corolario}
\label{cor:media}
Nas hipóteses do Teorema~\ref{thm:green-ret}, a circulação dividida pela área é a média de $\partial_x Q-\partial_y P$ em $D$. Se os retângulos encolhem até um ponto interior onde essa função é contínua, o quociente circulação por área tende ao rotacional escalar nesse ponto.
\end{corolario}

O rotacional escalar é a circulação local por unidade de área. Uma pá de rotor posta no ponto gira com a intensidade $\partial_x Q-\partial_y P$. Green é essa leitura infinitesimal somada no domínio inteiro.

\section{Regiões entre gráficos}

Um retângulo é o caso em que os gráficos são constantes. A mesma diferença de gráficos vale em regiões do tipo usual do cálculo.

\begin{definicao}
Uma região plana $D$ é do \emph{tipo~I} quando
\[
D=\bigl\{(x,y):a\le x\le b,\; g_1(x)\le y\le g_2(x)\bigr\},
\]
com $g_1$ e $g_2$ de classe $C^1$ em $[a,b]$ e $g_1\le g_2$. É do \emph{tipo~II} quando
\[
D=\bigl\{(x,y):c\le y\le d,\; h_1(y)\le x\le h_2(y)\bigr\},
\]
com $h_1$ e $h_2$ de classe $C^1$ em $[c,d]$ e $h_1\le h_2$.
\end{definicao}

\begin{proposicao}
\label{prop:tipoI}
Seja $D$ do tipo~I e seja $P$ de classe $C^1$ num aberto que contém $D$. Com $\partial D$ anti-horária,
\[
\oint_{\partial D} P\,\dif x=-\iint_D \partial_y P\,\dif A.
\]
\end{proposicao}

\begin{proof}
A fronteira anti-horária percorre o gráfico inferior $y=g_1(x)$ com $x$ de $a$ até $b$, depois o eventual segmento vertical em $x=b$, depois o gráfico superior $y=g_2(x)$ com $x$ de $b$ até $a$, depois o eventual segmento vertical em $x=a$. Nos segmentos verticais, $\dif x=0$, então eles não entram em $\oint P\,\dif x$. Resta
\begin{align*}
\oint_{\partial D} P\,\dif x
&=\int_a^b P\bigl(x,g_1(x)\bigr)\,\dif x+\int_b^a P\bigl(x,g_2(x)\bigr)\,\dif x\\
&=\int_a^b\Bigl(P\bigl(x,g_1(x)\bigr)-P\bigl(x,g_2(x)\bigr)\Bigr)\dif x.
\end{align*}
Para cada $x$, o teorema fundamental do cálculo entre os dois gráficos dá
\[
P\bigl(x,g_2(x)\bigr)-P\bigl(x,g_1(x)\bigr)
=\int_{g_1(x)}^{g_2(x)}\partial_y P(x,y)\,\dif y.
\]
O sinal troca e a integral dupla aparece:
\[
\oint_{\partial D} P\,\dif x
=-\int_a^b\int_{g_1(x)}^{g_2(x)}\partial_y P(x,y)\,\dif y\,\dif x
=-\iint_D\partial_y P\,\dif A.
\]
\end{proof}

\begin{proposicao}
\label{prop:tipoII}
Seja $D$ do tipo~II e seja $Q$ de classe $C^1$ num aberto que contém $D$. Com $\partial D$ anti-horária,
\[
\oint_{\partial D} Q\,\dif y=\iint_D \partial_x Q\,\dif A.
\]
\end{proposicao}

\begin{proof}
A fronteira anti-horária sobe pelo gráfico da direita $x=h_2(y)$, com $y$ de $c$ até $d$, e desce pelo gráfico da esquerda $x=h_1(y)$, com $y$ de $d$ até $c$. Os eventuais segmentos horizontais têm $\dif y=0$ e não entram em $\oint Q\,\dif y$. Portanto
\begin{align*}
\oint_{\partial D} Q\,\dif y
&=\int_c^d Q\bigl(h_2(y),y\bigr)\,\dif y+\int_d^c Q\bigl(h_1(y),y\bigr)\,\dif y\\
&=\int_c^d\Bigl(Q\bigl(h_2(y),y\bigr)-Q\bigl(h_1(y),y\bigr)\Bigr)\dif y\\
&=\int_c^d\int_{h_1(y)}^{h_2(y)}\partial_x Q(x,y)\,\dif x\,\dif y
=\iint_D\partial_x Q\,\dif A.
\end{align*}
A terceira igualdade é o teorema fundamental do cálculo na variável $x$, entre os dois gráficos.
\end{proof}

\begin{teorema}
\label{thm:green}
Seja $D$ uma região que é ao mesmo tempo do tipo~I e do tipo~II, e sejam $P$ e $Q$ de classe $C^1$ num aberto que contém $D$. Com $\partial D$ anti-horária,
\[
\oint_{\partial D} P\,\dif x+Q\,\dif y
=\iint_D\bigl(\partial_x Q-\partial_y P\bigr)\,\dif A.
\]
O mesmo vale se $D$ é união finita de regiões desse tipo, desde que as fronteiras internas sejam arcos $C^1$ por pedaços percorridos uma vez em cada sentido.
\end{teorema}

\begin{proof}
Nas regiões que admitem as duas descrições, a Proposição~\ref{prop:tipoI} controla $\oint P\,\dif x$ e a Proposição~\ref{prop:tipoII} controla $\oint Q\,\dif y$. A soma é a fórmula de Green.

Se $D=D_1\cup\cdots\cup D_m$, com interiores disjuntos e cada $D_k$ dos dois tipos, some as $m$ identidades. Num arco interno compartilhado por $D_k$ e $D_\ell$, a orientação anti-horária de $D_k$ deixa $D_k$ à esquerda de quem percorre o arco, e a orientação anti-horária de $D_\ell$ deixa $D_\ell$ à esquerda. Quem caminha com a região à esquerda percorre o arco compartilhado nos dois sentidos opostos. As integrais de $P\,\dif x+Q\,\dif y$ nesses dois percursos são opostas e se cancelam. O que não cancela é a fronteira externa de $D$, orientada com $D$ à esquerda, isto é, no sentido anti-horário da união.
\end{proof}

Dois retângulos lado a lado ilustram o cancelamento. O da esquerda sobe a aresta comum; o da direita, para manter o interior à esquerda, desce a mesma aresta. A soma das circulações esquece a parede interna.

\section{Área por Green}

\begin{proposicao}
\label{prop:area}
Seja $D$ uma região nas hipóteses do Teorema~\ref{thm:green}, com fronteira anti-horária. A área de $D$ admite as três leituras
\[
\mathrm{Area}(D)
=\oint_{\partial D}-y\,\dif x
=\oint_{\partial D}x\,\dif y
=\frac12\oint_{\partial D}-y\,\dif x+x\,\dif y.
\]
\end{proposicao}

\begin{proof}
A área é $\iint_D 1\,\dif A$. Basta escolher $P$ e $Q$ com $\partial_x Q-\partial_y P=1$ e aplicar Green.
\begin{itemize}
\item $P=0$, $Q=x$: então $\partial_x Q-\partial_y P=1$, e a área é $\oint x\,\dif y$.
\item $P=-y$, $Q=0$: então $\partial_x Q-\partial_y P=0-(-1)=1$, e a área é $\oint -y\,\dif x$.
\item $P=-y/2$, $Q=x/2$: então $\partial_x Q=1/2$ e $-\partial_y P=1/2$, e a área é $\frac12\oint -y\,\dif x+x\,\dif y$.
\end{itemize}
\end{proof}

\begin{exemplo}[Elipse]
A elipse $x=a\cos t$, $y=b\sin t$, com $a>0$, $b>0$ e $t\in[0,2\pi]$, é percorrida no sentido anti-horário. Use a terceira fórmula. Então $\dif x=-a\sin t\,\dif t$ e $\dif y=b\cos t\,\dif t$, de modo que
\begin{align*}
-y\,\dif x+x\,\dif y
&=-(b\sin t)(-a\sin t\,\dif t)+(a\cos t)(b\cos t\,\dif t)\\
&=ab(\sin^2 t+\cos^2 t)\,\dif t
=ab\,\dif t.
\end{align*}
A área é
\[
\frac12\int_0^{2\pi} ab\,\dif t=\pi ab.
\]
As outras duas fórmulas dão o mesmo número: $\int_0^{2\pi} x\,\dif y=\int_0^{2\pi} ab\cos^2 t\,\dif t=\pi ab$, e $\int_0^{2\pi}-y\,\dif x=\int_0^{2\pi} ab\sin^2 t\,\dif t=\pi ab$. A parametrização $(a\cos t,\,b\sin t)$ é $C^\infty$. Os gráficos $y=\pm b\sqrt{1-x^2/a^2}$ têm tangente vertical em $x=\pm a$, então a hipótese de gráficos $C^1$ até o bordo falha nesses dois pontos. A circulação acima não deriva esses gráficos. A identificação com a área é o caso $\partial_x Q-\partial_y P=1$: numa elipse um pouco menor, $x^2/a^2+y^2/b^2\le r^2$ com $r<1$, corte-se antes dos vértices por segmentos verticais; aí os gráficos são $C^1$, Green se aplica, e no limite $r\to 1$ os segmentos encolhem e a área tende a $\pi ab$.
\end{exemplo}

\section{Rotacional, divergência e o fluxo por volume}

\begin{definicao}
Seja $\vet{F}=(P,Q,R)$ de classe $C^1$ num aberto de $\R^3$. A \emph{divergência} e o \emph{rotacional} são
\begin{align*}
\nabla\cdot\vet{F}
&=\partial_x P+\partial_y Q+\partial_z R,\\
\nabla\times\vet{F}
&=\bigl(
\partial_y R-\partial_z Q,\;
\partial_z P-\partial_x R,\;
\partial_x Q-\partial_y P
\bigr).
\end{align*}
\end{definicao}

A divergência é a fonte local: fluido que nasce no ponto, por unidade de volume. O rotacional é a circulação local, a pá do rotor. As duas identidades seguintes dizem que um gradiente não tem pá, e que um rotor não tem fonte.

\begin{proposicao}
\label{prop:identidades}
Seja $\phi$ de classe $C^2$ e seja $\vet{F}$ de classe $C^2$ num aberto de $\R^3$. Então $\nabla\times(\nabla\phi)=\vet{0}$ e $\nabla\cdot(\nabla\times\vet{F})=0$.
\end{proposicao}

\begin{proof}
A hipótese $C^2$ dá parciais de segunda ordem contínuas, e Clairaut se aplica. Com $P=\partial_x\phi$, $Q=\partial_y\phi$ e $R=\partial_z\phi$,
\begin{align*}
\partial_y R-\partial_z Q&=\partial_y\partial_z\phi-\partial_z\partial_y\phi=0,\\
\partial_z P-\partial_x R&=\partial_z\partial_x\phi-\partial_x\partial_z\phi=0,\\
\partial_x Q-\partial_y P&=\partial_x\partial_y\phi-\partial_y\partial_x\phi=0.
\end{align*}
Para a divergência do rotacional, expanda
\begin{align*}
\nabla\cdot(\nabla\times\vet{F})
&=\partial_x(\partial_y R-\partial_z Q)
+\partial_y(\partial_z P-\partial_x R)
+\partial_z(\partial_x Q-\partial_y P)\\
&=\partial_x\partial_y R-\partial_x\partial_z Q
+\partial_y\partial_z P-\partial_y\partial_x R
+\partial_z\partial_x Q-\partial_z\partial_y P.
\end{align*}
Os pares $\partial_x\partial_y R-\partial_y\partial_x R$, $-\partial_x\partial_z Q+\partial_z\partial_x Q$ e $\partial_y\partial_z P-\partial_z\partial_y P$ se anulam por Clairaut.
\end{proof}

\begin{proposicao}
\label{prop:limite-div}
Seja $\vet{F}$ de classe $C^1$ num aberto que contém a origem. No cubo $Q_h=[0,h]^3$, com $h>0$ pequeno e com a fronteira orientada pela normal para fora, seja $\Phi(h)$ o fluxo de $\vet{F}$. Então
\[
\lim_{h\to 0^+}\frac{\Phi(h)}{h^3}=\nabla\cdot\vet{F}(\vet{0}).
\]
O mesmo limite, no cubo $\vet{x}_0+[0,h]^3$, recupera $\nabla\cdot\vet{F}(\vet{x}_0)$.
\end{proposicao}

\begin{proof}
Como $\vet{F}$ é $C^1$, é diferenciável na origem:
\[
\vet{F}(\vet{x})=\vet{F}(\vet{0})+A\vet{x}+\vet{R}(\vet{x}),
\]
onde $A=D\vet{F}(\vet{0})$ tem entradas $A_{ij}=\partial_j F_i(\vet{0})$, $\vet{R}(\vet{0})=\vet{0}$ e $\norm{\vet{R}(\vet{x})}/\norm{\vet{x}}\to 0$ quando $\vet{x}\to\vet{0}$. O fluxo é linear, então $\Phi=\Phi_{\mathrm{const}}+\Phi_{\mathrm{lin}}+\Phi_{\mathrm{resto}}$.

O campo constante $\vet{c}=\vet{F}(\vet{0})$ contribui $c_1 h^2$ na face $x=h$, cuja normal externa é $(1,0,0)$, e contribui $-c_1 h^2$ na face $x=0$, cuja normal externa é $(-1,0,0)$. A soma do par é zero. Os pares $y=0$, $y=h$ e $z=0$, $z=h$ se cancelam do mesmo modo. Logo $\Phi_{\mathrm{const}}=0$.

Na face de $x$ constante a normal externa é $(\pm 1,0,0)$, então só $F_1$ contribui. A parte linear $L_1=A_{11}x+A_{12}y+A_{13}z$ satisfaz
\[
L_1(h,y,z)-L_1(0,y,z)=A_{11}h.
\]
Os termos em $y$ e em $z$ aparecem nas duas faces com o mesmo valor e se cancelam no fluxo externo. Integrando no quadrado de área $h^2$,
\[
\iint_{[0,h]^2} A_{11}h\,\dif y\,\dif z=A_{11}h^3.
\]
O par de faces $y=\mathrm{constante}$ contribui $A_{22}h^3$, e o par $z=\mathrm{constante}$ contribui $A_{33}h^3$. Portanto
\[
\Phi_{\mathrm{lin}}=\bigl(A_{11}+A_{22}+A_{33}\bigr)h^3
=\bigl(\nabla\cdot\vet{F}(\vet{0})\bigr)\,h^3.
\]

Resta o erro. Dado $\eps>0$, existe $\delta>0$ tal que $0<\norm{\vet{x}}<\delta$ implica $\norm{\vet{R}(\vet{x})}\le\eps\norm{\vet{x}}$. No cubo, $\norm{\vet{x}}\le h\sqrt{3}$. Se $0<h<\delta/\sqrt{3}$, vale $\norm{\vet{R}(\vet{x})}\le\eps h\sqrt{3}$ em todo o cubo. A contribuição de $\vet{R}$ ao par de faces $x=\mathrm{constante}$ é
\[
\iint_{[0,h]^2}\bigl(R_1(h,y,z)-R_1(0,y,z)\bigr)\,\dif y\,\dif z,
\]
cujo módulo não passa de $2\cdot\eps h\sqrt{3}\cdot h^2=2\sqrt{3}\,\eps\, h^3$. Três pares de faces dão
\[
\bigl|\Phi_{\mathrm{resto}}(h)\bigr|\le 6\sqrt{3}\,\eps\, h^3.
\]
Juntando as três partes,
\[
\Biggl|\frac{\Phi(h)}{h^3}-\nabla\cdot\vet{F}(\vet{0})\Biggr|\le 6\sqrt{3}\,\eps.
\]
O número $\eps>0$ é arbitrário, logo o limite vale $\nabla\cdot\vet{F}(\vet{0})$. O volume do cubo é $h^3$, então o fluxo por unidade de volume tende à divergência. A translação $\vet{x}\mapsto\vet{x}-\vet{x}_0$ leva o cubo apoiado em $\vet{x}_0$ ao cubo apoiado na origem e repete a conta para $D\vet{F}(\vet{x}_0)$.
\end{proof}

\begin{intuicao}
Divergência positiva é fonte: o fluxo que sai de um cubo pequeno supera o que entra, e o excesso por unidade de volume é $\nabla\cdot\vet{F}$. Divergência negativa é sumidouro. Rotacional é o análogo para a circulação, já visível no Corolário~\ref{cor:media}. Green, Stokes e Gauss integram essa densidade e devolvem o total na fronteira.
\end{intuicao}

\section{Superfície paramétrica e fluxo}

\begin{definicao}
Seja $\vet{r}\colon D\subset\R^2\to\R^3$ de classe $C^1$, com $D$ limitado e de fronteira $C^1$ por pedaços. O \emph{produto fundamental} é $\vet{r}_u\times\vet{r}_v$. Onde ele não se anula, é normal à superfície $S=\vet{r}(D)$. A área é
\[
\mathrm{Area}(S)=\iint_D\norm{\vet{r}_u\times\vet{r}_v}\,\dif u\,\dif v,
\]
quando $\vet{r}$ é injetora no interior de $D$, salvo talvez numa união finita de curvas. Se $\vet{F}$ é contínuo num aberto que contém $S$, o \emph{fluxo} de $\vet{F}$ através de $S$, com a orientação de $\vet{r}_u\times\vet{r}_v$, é
\[
\iint_S\vet{F}\cdot\dif\vet{S}
=\iint_D\vet{F}\bigl(\vet{r}(u,v)\bigr)\cdot\bigl(\vet{r}_u\times\vet{r}_v\bigr)\,\dif u\,\dif v.
\]
Trocar a ordem dos parâmetros troca o sinal de $\vet{r}_u\times\vet{r}_v$ e o sinal do fluxo.
\end{definicao}

A orientação é a escolha de um dos dois sentidos da normal, de modo contínuo ao longo de $S$ quando $S$ é orientável. O vetor $\vet{r}_u\times\vet{r}_v$ escolhe um sentido. A regra da mão direita liga esse sentido ao bordo: os dedos acompanham $\partial S$ e o polegar aponta como a normal.

\begin{exemplo}[Hemisfério]
\label{ex:hemisferio}
Seja $S$ o hemisfério $x^2+y^2+z^2=1$, $z\ge 0$, e seja $\vet{F}=(0,0,z)$. Oriente $S$ pela normal de componente $z$ não negativa. A convenção esférica usada aqui é a do curso: $\varphi$ parte do eixo $z$ positivo e $\theta$ é o azimute no plano $xy$. A parametrização
\[
\vet{r}(\varphi,\theta)=(\sin\varphi\cos\theta,\;\sin\varphi\sin\theta,\;\cos\varphi),
\]
com $\varphi\in[0,\pi/2]$ e $\theta\in[0,2\pi]$, cobre $S$. As derivadas parciais
\begin{align*}
\vet{r}_\varphi&=(\cos\varphi\cos\theta,\;\cos\varphi\sin\theta,\;-\sin\varphi),\\
\vet{r}_\theta&=(-\sin\varphi\sin\theta,\;\sin\varphi\cos\theta,\;0)
\end{align*}
são contínuas no retângulo fechado dos parâmetros. Com $\vet{a}=\vet{r}_\varphi$ e $\vet{b}=\vet{r}_\theta$, as componentes de $\vet{a}\times\vet{b}$ são
\begin{align*}
a_2 b_3-a_3 b_2
&=\sin\varphi\cdot\sin\varphi\cos\theta
=\sin^2\varphi\cos\theta,\\
a_3 b_1-a_1 b_3
&=(-\sin\varphi)(-\sin\varphi\sin\theta)
=\sin^2\varphi\sin\theta,\\
a_1 b_2-a_2 b_1
&=\cos\varphi\sin\varphi\cos^2\theta+\cos\varphi\sin\varphi\sin^2\theta
=\sin\varphi\cos\varphi.
\end{align*}
Logo
\[
\vet{r}_\varphi\times\vet{r}_\theta
=\bigl(\sin^2\varphi\cos\theta,\;\sin^2\varphi\sin\theta,\;\sin\varphi\cos\varphi\bigr)
=\sin\varphi\;\vet{r}(\varphi,\theta).
\]
No polo $\varphi=0$ esse produto se anula. O integrando que aparece abaixo é contínuo, e um ponto não altera a integral.
A terceira componente $\sin\varphi\cos\varphi$ é $\ge 0$ em $[0,\pi/2]$, e é positiva no hemisfério aberto. Esta é a orientação pedida. Como $\vet{F}(\vet{r})=(0,0,\cos\varphi)$,
\[
\vet{F}(\vet{r})\cdot(\vet{r}_\varphi\times\vet{r}_\theta)
=\cos\varphi\cdot\sin\varphi\cos\varphi
=\sin\varphi\cos^2\varphi.
\]
O fluxo é
\[
\int_0^{2\pi}\int_0^{\pi/2}\sin\varphi\cos^2\varphi\,\dif\varphi\,\dif\theta
=2\pi\Biggl[-\frac{\cos^3\varphi}{3}\Biggr]_0^{\pi/2}
=2\pi\Biggl(0-\Bigl(-\frac13\Bigr)\Biggr)
=\frac{2\pi}{3}.
\]
O campo aponta para cima no semiespaço $z>0$ e a normal escolhida também tem componente $z$ positiva, então o fluxo positivo é o sinal esperado. O valor $2\pi/3$ é menor que a área $2\pi$ do hemisfério: a componente normal de $\vet{F}$ não vale $1$ o tempo todo.
\end{exemplo}

\section{Stokes}

\begin{teorema}
\label{thm:stokes}
Seja $S$ uma superfície orientável, imagem de uma parametrização $\vet{r}$ de classe $C^2$ definida numa região plana $D$ à qual se aplica o Teorema~\ref{thm:green}, com $\vet{r}_u\times\vet{r}_v$ nunca zero no interior e apontando no sentido da orientação de $S$. Seja $\partial S$ a curva fechada simples imagem de $\partial D$, orientada pela regra da mão direita. Seja $\vet{F}$ de classe $C^1$ num aberto que contém $S$. Então
\[
\oint_{\partial S}\vet{F}\cdot\dif\vet{r}
=\iint_S(\nabla\times\vet{F})\cdot\dif\vet{S}.
\]
\end{teorema}

\begin{corolario}
\label{cor:green-stokes}
Green é Stokes com $S$ contida no plano $xy$. Se $\vet{F}=(P,Q,R)$ e $S=D$ tem normal $\vet{k}$, então $(\nabla\times\vet{F})\cdot\vet{k}=\partial_x Q-\partial_y P$ e, no bordo plano, $\dif z=0$, de modo que a circulação é $\oint P\,\dif x+Q\,\dif y$.
\end{corolario}

\begin{proof}
A parametrização $\vet{r}(x,y)=(x,y,0)$ tem $\vet{r}_x\times\vet{r}_y=(0,0,1)$. A terceira componente do rotacional é $\partial_x Q-\partial_y P$. O Teorema~\ref{thm:stokes} devolve exatamente a fórmula de Green, com o sentido anti-horário correspondente ao polegar em $\vet{k}$.
\end{proof}

\begin{esqueleto}
\item A parametrização $\vet{r}(u,v)$, $(u,v)\in D$, é escolhida de modo que $\vet{r}_u\times\vet{r}_v$ aponte como a normal de $S$, e $\partial D$ anti-horário seja levado no bordo $\partial S$.
\item Se $\gamma(t)=(u(t),v(t))$ percorre $\partial D$, a curva no espaço é $\vet{r}(\gamma(t))$ e a sua derivada é $\vet{r}_u u'+\vet{r}_v v'$. Então
\[
\vet{F}(\vet{r})\cdot(\vet{r}_u u'+\vet{r}_v v')
=\tilde P\,u'+\tilde Q\,v',
\]
com $\tilde P(u,v)=\vet{F}(\vet{r})\cdot\vet{r}_u$ e $\tilde Q(u,v)=\vet{F}(\vet{r})\cdot\vet{r}_v$. Integrando,
\[
\oint_{\partial S}\vet{F}\cdot\dif\vet{r}
=\oint_{\partial D}\tilde P\,\dif u+\tilde Q\,\dif v.
\]
\item Como $\vet{F}$ é $C^1$ e $\vet{r}$ é $C^2$, os coeficientes $\tilde P$ e $\tilde Q$ são $C^1$ em $D$. Green no plano dos parâmetros dá
\[
\oint_{\partial D}\tilde P\,\dif u+\tilde Q\,\dif v
=\iint_D\bigl(\partial_u\tilde Q-\partial_v\tilde P\bigr)\,\dif u\,\dif v.
\]
\item Derive o produto. A regra do produto e a cadeia produzem
\begin{align*}
\partial_u\tilde Q
&=\bigl(D\vet{F}(\vet{r})\,\vet{r}_u\bigr)\cdot\vet{r}_v
+\vet{F}(\vet{r})\cdot\vet{r}_{vu},\\
\partial_v\tilde P
&=\bigl(D\vet{F}(\vet{r})\,\vet{r}_v\bigr)\cdot\vet{r}_u
+\vet{F}(\vet{r})\cdot\vet{r}_{uv}.
\end{align*}
As derivadas mistas de $\vet{r}$ coincidem, e os dois termos com $\vet{F}$ se cancelam. Sobra
\[
\partial_u\tilde Q-\partial_v\tilde P
=\sum_{i,j}\partial_j F_i
\bigl(\partial_u x_j\,\partial_v x_i-\partial_v x_j\,\partial_u x_i\bigr),
\]
onde $(x_1,x_2,x_3)=(x,y,z)$.
\item Passo da componente. O coeficiente de $\partial_x F_2$ é o par $(i,j)=(2,1)$:
\[
\partial_x F_2\,(x_u y_v-x_v y_u).
\]
O coeficiente de $\partial_y F_1$ é o par $(i,j)=(1,2)$:
\[
\partial_y F_1\,(y_u x_v-y_v x_u)
=-\partial_y F_1\,(x_u y_v-x_v y_u).
\]
A soma é $(\partial_x F_2-\partial_y F_1)(x_u y_v-x_v y_u)$, que é a terceira componente de $(\nabla\times\vet{F})\cdot(\vet{r}_u\times\vet{r}_v)$, porque a terceira componente de $\vet{r}_u\times\vet{r}_v$ é $x_u y_v-y_u x_v$. Os outros pares repetem o mesmo agrupamento:
\begin{align*}
(\partial_y F_3-\partial_z F_2)(y_u z_v-y_v z_u),\\
(\partial_z F_1-\partial_x F_3)(z_u x_v-x_u z_v).
\end{align*}
São a primeira e a segunda componentes do mesmo produto escalar. Portanto
\[
\partial_u\tilde Q-\partial_v\tilde P
=(\nabla\times\vet{F})\cdot(\vet{r}_u\times\vet{r}_v).
\]
\item A integral dupla em $D$ é, pela definição de fluxo, $\iint_S(\nabla\times\vet{F})\cdot\dif\vet{S}$.
\end{esqueleto}

O esqueleto reduz Stokes a Green no plano dos parâmetros. O passo que não é metáfora é o agrupamento das componentes: cada par $\partial_j F_i-\partial_i F_j$ multiplica a componente correspondente de $\vet{r}_u\times\vet{r}_v$.

\section{Gauss}

\begin{teorema}
\label{thm:gauss}
Seja $E$ um sólido elementar: união finita de regiões da forma
\[
\bigl\{(x,y,z):(x,y)\in D,\; g_1(x,y)\le z\le g_2(x,y)\bigr\},
\]
e também união finita de regiões análogas entre gráficos de $x$ e entre gráficos de $y$, com os gráficos de classe $C^1$ e $D$ nas hipóteses de Green. Seja $\vet{F}$ de classe $C^1$ num aberto que contém $E$. Oriente $\partial E$ pela normal para fora. Então
\[
\iint_{\partial E}\vet{F}\cdot\dif\vet{S}
=\iiint_E(\nabla\cdot\vet{F})\,\dif V.
\]
\end{teorema}

\begin{proof}
Basta provar o caso de uma região entre dois gráficos e de um campo vertical $\vet{F}=(0,0,R)$, e depois indicar as outras componentes. Seja
\[
E=\bigl\{(x,y,z):(x,y)\in D,\; g_1(x,y)\le z\le g_2(x,y)\bigr\},
\]
com $g_1\le g_2$ de classe $C^1$ num aberto que contém o compacto $D$. A fronteira tem três partes: o topo $z=g_2$, a base $z=g_1$ e a parede lateral sobre $\partial D$, vertical.

No topo, $\vet{r}(x,y)=(x,y,g_2(x,y))$ tem
\[
\vet{r}_x\times\vet{r}_y=(-\partial_x g_2,\,-\partial_y g_2,\,1),
\]
normal de componente $z$ positiva, que aponta para fora. Como $\vet{F}=(0,0,R)$,
\[
\vet{F}\cdot(\vet{r}_x\times\vet{r}_y)=R\bigl(x,y,g_2(x,y)\bigr).
\]
Na base, a normal para fora aponta para baixo. O produto $\vet{r}_x\times\vet{r}_y$ calculado com $g_1$ tem terceira componente $+1$ e aponta para dentro, então a orientação externa troca o sinal:
\[
\vet{F}\cdot\dif\vet{S}=-R\bigl(x,y,g_1(x,y)\bigr)\,\dif x\,\dif y.
\]
Na parede lateral a normal é horizontal, porque a parede é vertical, e $\vet{F}$ é vertical, então o produto escalar é zero.

O fluxo total é
\begin{align*}
\iint_D\Bigl(R\bigl(x,y,g_2\bigr)-R\bigl(x,y,g_1\bigr)\Bigr)\dif x\,\dif y
&=\iint_D\int_{g_1(x,y)}^{g_2(x,y)}\partial_z R(x,y,z)\,\dif z\,\dif x\,\dif y\\
&=\iiint_E\partial_z R\,\dif V
=\iiint_E(\nabla\cdot\vet{F})\,\dif V.
\end{align*}
A passagem da diferença $R(g_2)-R(g_1)$ para a integral de $\partial_z R$ é o teorema fundamental do cálculo.

Para $\vet{F}=(P,0,0)$ a mesma conta vale numa descrição de $E$ entre gráficos $x=h_1(y,z)$ e $x=h_2(y,z)$: o fluxo externo das duas tampas é $\iint\bigl(P(h_2)-P(h_1)\bigr)$ e isso é $\iiint\partial_x P$. Para $\vet{F}=(0,Q,0)$, usam-se gráficos em $y$. Um campo geral é a soma das três componentes.

Se $E$ é união finita de sólidos desse tipo, some as identidades. Cada face interna é externa para um pedaço e interna para o vizinho: as normais ``para fora de cada pedaço'' são opostas, e os fluxos se cancelam. Sobra a fronteira externa de $E$, orientada para fora.
\end{proof}

\section{Duas consequências e dois exemplos}

\begin{proposicao}
\label{prop:fluxo-rot}
Seja $\vet{F}$ de classe $C^2$ num aberto que contém um sólido elementar $E$. Então o fluxo de $\nabla\times\vet{F}$ através de $\partial E$, com normal para fora, é zero.
\end{proposicao}

\begin{proof}
Pela Proposição~\ref{prop:identidades}, $\nabla\cdot(\nabla\times\vet{F})=0$. O campo $\nabla\times\vet{F}$ é de classe $C^1$. O Teorema~\ref{thm:gauss} dá
\[
\iint_{\partial E}(\nabla\times\vet{F})\cdot\dif\vet{S}
=\iiint_E\nabla\cdot(\nabla\times\vet{F})\,\dif V=0.
\]
\end{proof}

\begin{proposicao}
\label{prop:circ-grad}
Seja $\phi$ de classe $C^2$ num aberto que contém uma superfície $S$ nas hipóteses do Teorema~\ref{thm:stokes}. Então a circulação de $\nabla\phi$ ao longo de $\partial S$ é zero.
\end{proposicao}

\begin{proof}
O gradiente $\nabla\phi$ é de classe $C^1$, e $\nabla\times(\nabla\phi)=\vet{0}$ pela Proposição~\ref{prop:identidades}. Stokes dá
\[
\oint_{\partial S}\nabla\phi\cdot\dif\vet{r}
=\iint_S\bigl(\nabla\times(\nabla\phi)\bigr)\cdot\dif\vet{S}=0.
\]
\end{proof}

A segunda proposição é o teorema do gradiente lido por Stokes, quando a curva fechada é bordo de alguma superfície dentro do domínio. No plano furado, a circunferência unitária não é bordo de uma superfície contida no domínio da forma angular: qualquer disco que ela limita contém a origem, onde o campo não é $C^1$.

\begin{exemplo}[Stokes mais curto que a linha]
\label{ex:stokes-disco}
Seja $\vet{F}=(z,x,y)$ e seja $C$ a circunferência $x^2+y^2=1$, $z=0$, anti-horária. O disco $S$ dado por $x^2+y^2\le 1$, $z=0$, tem normal $\vet{k}$ compatível com a regra da mão direita. O rotacional é
\begin{align*}
\nabla\times\vet{F}
&=\bigl(\partial_y(y)-\partial_z(x),\;
\partial_z(z)-\partial_x(y),\;
\partial_x(x)-\partial_y(z)\bigr)
=(1,1,1).
\end{align*}
Stokes lê só a terceira componente no disco:
\[
\oint_C\vet{F}\cdot\dif\vet{r}
=\iint_S 1\,\dif A=\pi.
\]
A integral de linha confirma, e é mais longa. Com $\vet{r}(t)=(\cos t,\sin t,0)$, $t\in[0,2\pi]$,
\[
\vet{F}(\vet{r}(t))=(0,\cos t,\sin t),\qquad
\vet{r}'(t)=(-\sin t,\cos t,0),
\]
e o produto escalar vale $\cos^2 t$. Então
\[
\int_0^{2\pi}\cos^2 t\,\dif t
=\int_0^{2\pi}\frac{1+\cos 2t}{2}\,\dif t
=\pi.
\]
Os dois caminhos dão $\pi$. O rotacional constante torna a superfície a conta curta.
\end{exemplo}

\begin{exemplo}[Gauss mais curto que o fluxo direto]
\label{ex:gauss-esfera}
Seja $\vet{F}=(x,y,z)$ e seja $S$ a esfera $x^2+y^2+z^2=R^2$, $R>0$, com normal para fora. A divergência é $\nabla\cdot\vet{F}=3$. O volume da bola é calculado pelo empilhamento de discos. A faixa entre $z$ e $z+\dif z$ tem raio $\sqrt{R^2-z^2}$ e volume de fatia $\pi(R^2-z^2)\,\dif z$, logo
\[
\mathrm{Vol}
=\int_{-R}^{R}\pi(R^2-z^2)\,\dif z
=\pi\Biggl[R^2 z-\frac{z^3}{3}\Biggr]_{-R}^{R}
=\pi\Biggl(\Bigl(R^3-\frac{R^3}{3}\Bigr)-\Bigl(-R^3+\frac{R^3}{3}\Bigr)\Biggr)
=\frac{4\pi R^3}{3}.
\]
O Teorema~\ref{thm:gauss}, aplicado à bola como sólido elementar, dá o fluxo
\[
\iint_S\vet{F}\cdot\dif\vet{S}
=\iiint_E 3\,\dif V
=3\cdot\frac{4\pi R^3}{3}
=4\pi R^3.
\]
A conta direta chega ao mesmo valor. Na parametrização $\vet{r}(\varphi,\theta)=R(\sin\varphi\cos\theta,\sin\varphi\sin\theta,\cos\varphi)$, com $\varphi\in[0,\pi]$ e $\theta\in[0,2\pi]$, cada derivada parcial traz um fator $R$ em relação ao hemisfério unitário, e o produto vetorial traz $R^2$. O cálculo de componentes feito no Exemplo~\ref{ex:hemisferio} dá, para raio $1$, o vetor $\sin\varphi\,\hat{\vet{r}}$. Aqui
\[
\vet{r}_\varphi\times\vet{r}_\theta
=R^2\sin\varphi\,\hat{\vet{r}}
=R\sin\varphi\,\vet{r},
\]
onde $\hat{\vet{r}}=\vet{r}/R$ e $\vet{F}(\vet{r})=\vet{r}$. Portanto
\[
\vet{F}\cdot(\vet{r}_\varphi\times\vet{r}_\theta)
=\vet{r}\cdot(R\sin\varphi\,\vet{r})
=R\sin\varphi\cdot R^2
=R^3\sin\varphi.
\]
Integrando,
\[
\int_0^{2\pi}\int_0^{\pi} R^3\sin\varphi\,\dif\varphi\,\dif\theta
=2\pi R^3\bigl[-\cos\varphi\bigr]_0^{\pi}
=2\pi R^3\bigl(1-(-1)\bigr)
=4\pi R^3.
\]
Gauss e o fluxo direto concordam. A esfera, vista como gráfico $z=\pm\sqrt{R^2-x^2-y^2}$ sobre o disco fechado de raio $R$, tem plano tangente vertical no equador, então esse par de gráficos não é $C^1$ até o bordo. Para usar o Teorema~\ref{thm:gauss} dentro da hipótese escrita, fixe $a\in(0,R)$ e seja $E_a$ a parte da bola sobre o disco $x^2+y^2\le a^2$. Aí a função $\sqrt{R^2-x^2-y^2}$ é $C^1$ numa vizinhança do disco, porque $R^2-a^2>0$. A fronteira de $E_a$ é formada pelas duas calotas e pelo cilindro $x^2+y^2=a^2$. Nesse cilindro a normal é horizontal, e a componente $(0,0,z)$ tem fluxo zero. O teorema dá que o fluxo de $(0,0,z)$ pelas calotas é o volume de $E_a$. A faixa esférica que falta corresponde a $\varphi\in[\varphi_a,\pi-\varphi_a]$, com $\sin\varphi_a=a/R$. A área dela é
\[
\int_0^{2\pi}\int_{\varphi_a}^{\pi-\varphi_a} R^2\sin\varphi\,\dif\varphi\,\dif\theta
=2\pi R^2\cdot 2\cos\varphi_a
=4\pi R\sqrt{R^2-a^2},
\]
usando $\cos\varphi_a=\sqrt{R^2-a^2}/R$ e o elemento de área $R^2\sin\varphi$ já obtido no produto fundamental. Essa área tende a zero quando $a\to R$, e $\vet{F}$ é limitado na esfera, então o fluxo pela faixa tende a zero. As componentes $(x,0,0)$ e $(0,y,0)$ repetem o argumento nos outros eixos, e cada uma contribui o volume $4\pi R^3/3$. A soma dos três fluxos é de novo $4\pi R^3$.
\end{exemplo}

\section{Exercícios}

\begin{exercicio}[{[Básico --- Green num retângulo]}]
Seja $D=[0,1]\times[0,2]$, com fronteira anti-horária, $P=-y$ e $Q=x$. Calcule $\partial_x Q-\partial_y P$ e a integral dupla. Calcule também a circulação pelos quatro lados.
\end{exercicio}

\begin{exercicio}[{[Básico --- divergência e rotacional]}]
Para $\vet{F}=(xy,\, yz,\, zx)$, calcule $\nabla\cdot\vet{F}$ e $\nabla\times\vet{F}$.
\end{exercicio}

\begin{exercicio}[{[Básico --- área do círculo]}]
Use $\frac12\oint -y\,\dif x+x\,\dif y$ para provar que o disco $x^2+y^2\le R^2$, $R>0$, tem área $\pi R^2$. A circunferência é $x=R\cos t$, $y=R\sin t$, $t\in[0,2\pi]$.
\end{exercicio}

\begin{exercicio}[{[Básico --- área e fluxo de um gráfico]}]
Seja $\vet{r}(u,v)=(u,v,u+v)$ com $(u,v)\in[0,1]^2$, e seja $\vet{F}=(0,0,1)$. Calcule $\vet{r}_u\times\vet{r}_v$, a área da superfície e o fluxo de $\vet{F}$.
\end{exercicio}

\begin{exercicio}[{[Intermediário --- circulação pelo Green]}]
Seja $C$ a fronteira anti-horária do quadrado $[0,1]^2$. Calcule
\[
\oint_C -y^2\,\dif x+x\,\dif y
\]
pelo teorema de Green. Registre a hipótese de classe $C^1$ que autoriza o uso.
\end{exercicio}

\begin{exercicio}[{[Intermediário --- Stokes no disco]}]
Seja $\vet{F}=(z,x,y)$ e seja $C$ a circunferência unitária do plano $z=0$, anti-horária. Calcule a circulação pelo teorema de Stokes, usando o disco como superfície. Confira pela integral de linha.
\end{exercicio}

\begin{exercicio}[{[Intermediário --- Gauss no cubo]}]
Seja $E=[0,1]^3$ e seja $\vet{F}=(x,y,z)$, com a fronteira orientada para fora. Calcule o fluxo pelo teorema da divergência. Confira somando o fluxo das seis faces.
\end{exercicio}

\begin{exercicio}[{[Avançado --- Green em região tipo II]}]
Seja $D$ do tipo~II, com $h_1$ e $h_2$ de classe $C^1$, e seja $Q$ de classe $C^1$ num aberto que contém $D$. Prove
\[
\oint_{\partial D} Q\,\dif y=\iint_D\partial_x Q\,\dif A
\]
com a fronteira anti-horária, escrevendo os dois gráficos e o uso do teorema fundamental do cálculo. Não diga apenas que a conta é análoga à do retângulo.
\end{exercicio}

\begin{exercicio}[{[Avançado --- Green não vê o furo]}]
Seja
\[
\vet{F}=\Biggl(-\frac{y}{x^2+y^2},\;\frac{x}{x^2+y^2}\Biggr)
\]
em $\R^2\setminus\{\vet{0}\}$, e seja $C$ a circunferência unitária anti-horária.
\begin{enumerate}
\item Calcule $\partial_x Q-\partial_y P$ em todo ponto do domínio.
\item Calcule $\oint_C\vet{F}\cdot\dif\vet{r}$.
\item A integral dupla formal de $\partial_x Q-\partial_y P$ no disco unitário seria $0$. Explique qual hipótese do teorema de Green falha, e por que os dois números podem ser diferentes.
\end{enumerate}
\end{exercicio}

\end{document}
